\section{Android Repackaging and Artifact-Level Evidence}
\label{sec:repackaging}

Repackaging begins with an existing application, a modification, re-signing and redistribution. The modification may change advertising, branding, payment behavior or malicious functionality. A useful analysis pipeline separates candidate retrieval, region comparison, pair classification, historical direction, behavior and release authorization. RePack's 57-approach review makes the importance of these distinct evaluation contracts explicit~\cite{repack}. We organize the mechanisms by the information each preserves and the operational decision it enables.

\subsection{Code representations and candidate retrieval}

Normalized fingerprints offer compact comparison when names or local syntax change; richer structural representations preserve relations at greater extraction cost. DroidMOSS uses fuzzy code signatures, DNADroid compares method data-dependence graphs after metadata retrieval, and AnDarwin uses successive approximate clusterings~\cite{droidmoss,attackclones,andarwin}. Juxtapp compares code-feature sets; WuKong filters candidates before detailed comparison~\cite{juxtapp,wukong}. This staged design avoids spending expensive confirmation work on every application pair, but a first-stage miss cannot be repaired by a more precise second stage.

The feature must match the relation sought. Declaration-only class signatures reveal shared interfaces, while normalized method structures and graph-derived features retain different implementation information~\cite{reuseandroid,codematch,simidroid}. Normalization improves resilience to selected transformations and also merges some distinct implementations. Libraries, templates, same-developer families and independently produced compatible code are therefore important alternatives to unauthorized copying. The early plagiarism study's synthetic derivatives and resampled scalability sets supply a different oracle from naturally occurring application histories~\cite{plagiarizing}.

RepDetector's bounded symbolic-output comparisons offer a semantic view for selected code, not whole-program equivalence; loops and native calls remain scoped~\cite{semanticsrepack}. CLANdroid instead retrieves functionally related apps, not labeled repackages~\cite{similarapps}. Its differing recommendation and participant outcomes illustrate why a retrieval ranking and a copying verdict cannot share one unqualified accuracy label.

Market-scale systems combine indexing, filtering and confirmation~\cite{accuracyclone,wukong,massvet}; APPraiser explicitly retains multiple relatedness categories and an unknown state~\cite{appraiser}. Their engineering lesson is positive: select a candidate architecture for the expected review workload, then report its extraction, retrieval and confirmation costs separately. AnDarwin's online full-app path, the centroid system's incremental path and Juxtapp's new--new/new--old updates have different prerequisites~\cite{andarwin,accuracyclone,juxtapp}. WuKong's per-submission update path is proposed rather than measured~\cite{wukong}. Offline feature construction, retained databases and human adjudication cannot be inferred from a fast lookup figure.

\subsection{Localizing changes and distinguishing behavior}

Piggybacking analysis asks where a mostly retained host differs from a candidate, so localized evidence can be more informative than a single application score~\cite{piggybacking}. DR-Droid partitions recovered class relations; DAPASA focuses on a selected sensitive subgraph~\cite{heterogeneity,dapasa}. The former's operational partition drops direction although its stated graph definition uses directed reachability; neither this partition nor a static sensitive subgraph supplies the insertion history or every Android runtime dependency.

A stronger claim joins a plausible host relationship to evidence that the differing region is an added transformation rather than ordinary product variation. Controlled original--derivative construction, deployed marks or authenticated source history can supply that link. Behavioral analysis then evaluates what the changed code does. DREBIN and MaMaDroid classify static features or behavioral abstractions, while FlowDroid analyzes lifecycle-aware information flow~\cite{drebin,mamadroid,flowdroid}. These support property-specific \claimlevel{6} evidence, not ancestry by themselves.

The 2019 MaMaDroid extension evaluates additional abstractions and a frequency variant; it is not interchangeable with the conference paper~\cite{mamadroid}. Static sequences are not executed traces, and classification latency excludes extraction and maintenance. DREBIN likewise explains feature-score contributions rather than runtime events~\cite{drebin}. Label construction, unsupported inputs and preprocessing remain part of each comparison contract. A localized anomalous region can be a valuable investigation lead without being a proved copying event.

\subsection{Resources, interfaces and designed origin evidence}

Resource sets, layout structure, rendered interfaces and navigation graphs preserve presentation information that code-only views can miss~\cite{resourcedriven,resourceeval,viewdroid}. They are particularly useful for brand impersonation and candidates with substantial code rewriting. Localization, restyling, template reuse and sampled screen states define their scope. DroidEagle combines trusted-reference layout comparison with certificate fingerprints~\cite{droideagle}; this supplies more than certificate presence alone, but a different key still requires external authorization evidence.

Resources also expose common-code confounding. FSquaDRA's largest two returned clusters were judged legitimate because ActionBarSherlock dominated the shared files~\cite{fsquadra}. Agreement with another similarity tool did not resolve history. ResDroid and ViewDroid use different template-positive oracles, and Juxtapp separates common-builder and same-author reuse~\cite{resourcedriven,viewdroid,juxtapp}. Combining code and UI is useful when they challenge different alternatives; correlated generated bindings should not be counted as independent corroboration.

Embedded marks make another design choice: add origin evidence before a dispute rather than infer it only from the resulting artifact. AppInk recovers a developer-chosen runtime graph using a confidential companion manifest~\cite{appink}. Its evaluated transformations, three-app timing study and trusted-arbitrator assumptions remain bounded; source-reported repair of a renaming failure is not an independent rerun here. Prior deployment and claimant identity can support designed origin. A contractor may legitimately possess marked source, so origin remains distinct from policy-scoped release authority.

\subsection{Transformation and oracle contracts}

Ordinary builds and adversarial edits both change observability. Compiler transformations, shrinking, identifier changes, control-flow rewriting, packing, resource changes and signing must be described at the stage where they act. DroidChameleon's selected historical transformations test a security-relevant preservation contract, not all observable semantics or current products~\cite{droidchameleon}. CodeMatch/SimiDroid and resource-based methods occupy different transformation cells~\cite{codematch,simidroid,resourceeval}. Reference-based transformed-copy detection also differs from open-world discovery with an unknown ancestor.

\begin{table}[t]
\caption{Representative Android artifact-evidence families. The claim ceiling assumes no auxiliary provenance or policy evidence.}
\label{tab:android-families}
\small
\begin{tabularx}{\textwidth}{L{2.45cm}Y L{1.35cm}Y}
\toprule
Family and examples & Observation and modeled resilience & Ceiling & Recurring failure boundary \\
\midrule
Fuzzy and normalized code~\cite{droidmoss,codematch} & Method or application fingerprints; often tolerates renaming and local edits & \claimlevel{2} & Common libraries, templates, coarse collisions, unseen rewrites \\
Graph and structural abstraction~\cite{attackclones,andarwin} & Method data-dependence graphs or approximate graph-derived features & \claimlevel{2}; \claimlevel{3} with external pair oracle & Graph rewriting, incomplete recovery, direction supplied by labels \\
Market-scale matching~\cite{massvet,wukong,accuracyclone} & Indexed features and system-specific filtering; MassVet also raises conditional risk alerts & \claimlevel{2}; \claimlevel{6} in MassVet & Base-rate effects, sparse labels, candidate-stage false negatives \\
Resource sets~\cite{resourcedriven,resourceeval}; UI views~\cite{viewdroid,droideagle} & Resources, layouts, images, static navigation graphs, or rendered views & \claimlevel{2} & Templates, localization, state coverage, deliberate restyling \\
Embedded marks~\cite{appink} & Designed watermark verified in a candidate & \claimlevel{3} with deployment record & Removal, forgery, licensed reuse, weak identity binding \\
Malware classification~\cite{drebin,mamadroid,heterogeneity}; flow/subgraph evidence~\cite{flowdroid,dapasa} & Static features, sensitive subgraphs, flows, anomalous regions & \claimlevel{6} for modeled property & Does not identify ancestor; analysis unsoundness or incomplete environment \\
\bottomrule
\end{tabularx}
\end{table}

A hard-negative set should include legitimate versions, shared libraries, common builders, authorized variants and independently generated siblings. Random unrelated pairs are usually easier. Family, developer, market, time and prompt-family boundaries matter when partitioning datasets; otherwise a database or threshold can reveal information unavailable at the intended decision time.

Direction requires an oracle beyond similarity. DNADroid admits a common predecessor, DroidMOSS assumes an official-market original, and AdRob distinguishes certificate identity from owner identity~\cite{attackclones,droidmoss,adrob}. WuKong and APPraiser retain chronology/reference-presence conditions; a nominal piggybacking host may itself be repackaged~\cite{wukong,appraiser,piggybacking}. DroidKin's public debug certificates show why certificate text need not name one developer~\cite{droidkin}. These findings motivate explicit parent/transformation evidence rather than a universal market-order heuristic.

Evaluation units must also remain distinct. RePack separates known-positive recall from returned-candidate review; detector comparisons with extraction failures or different completed-pair denominators cannot be pooled~\cite{repack,repackeval,simidroid}. AndroZoo supplies access to artifact snapshots and scanner observations, not authenticated copy histories or validated scanner truth~\cite{androzoo}. Scanner thresholds, missing artifacts, transformation failures and unresolved labels should therefore remain visible. No original detector or transformation was executed for this review. Source-specific denominators, timing stages and representation locators remain in the companion evidence records.

\subsection{Selecting and composing artifact evidence}

For inexpensive triage, use indexed normalized features or resources matched to the expected transformation family. For disputed regions, retain module-aware structural or bounded semantic evidence, including unknown outcomes. For altered behavior, apply a property-specific analysis to the exact candidate. For directional or authorization decisions, add a historical or policy record rather than interpreting a stronger similarity score as a different kind of evidence.

An operational report should bind its output to the artifact, compared regions, ignored components, feature/database state, threshold, candidate universe and observed transformations. Retrieval curves, returned-candidate precision and known-positive recall describe different stages. A reviewable pipeline exposes each stage and its failure boundary, allowing artifact evidence to remain useful when trustworthy process records are absent.
